\documentclass[12pt,usepdftitle=false,aspectratio=1610]{beamer}

\usetheme{Boadilla}
%\usetheme{Madrid}
%\usecolortheme{seahorse}

\setbeamertemplate{navigation symbols}{}

%\usefonttheme[onlymath]{serif}
\usefonttheme[]{serif}
%\useinnertheme{circles}
\setbeamertemplate{caption}[numbered]

\usepackage[T2A]{fontenc}			
\usepackage[utf8]{inputenc}			
\usepackage[english,russian]{babel}
\usepackage{textcomp}
\usepackage{gensymb}
\usepackage{mathtext,cmap,multirow,amsmath,tcolorbox} 				
\usepackage{graphicx,wrapfig,subfig,mathtools,gensymb}
\usepackage[font=small,labelfont=bf]{caption}
\usepackage{minted}

\DeclarePairedDelimiter{\norm}{\lVert}{\rVert} 

\graphicspath{{./figs/08}}

\title[Лекция 8]{Проективная модель камеры}

\author{Александр Танченко}
\institute{}
\date{2025}

\begin{document}

%------------------------------------------------------------
\begin{frame}
    \titlepage
\end{frame}

%------------------------------------------------------------
\begin{frame}
    \frametitle{Центральная проекция в евклидовых координатах}
    \begin{figure}
        \centering
        \includegraphics[height=0.3\textheight]{fig1.pdf}
        \includegraphics[height=0.3\textheight]{fig2.pdf}
    \end{figure}
    Pinhole камера:
    \begin{itemize}
        \item \textbf{image plane} -- фокальная плоскость
        \item $\mathbf{p}$ -- оптический центр  камеры
        \item $f=\norm{\mathbf{p}-\mathbf{C}}$ -- фокусное расстояние
        \item \textbf{principal axis} -- оптическая ось 
        \item \textbf{principal plane} -- плоскость, параллельная image plane и содержащая $\mathbf{C}$
    \end{itemize}
    \medskip

    \textbf{Центральная проекция} (из подобия треугольников):
    $$
        (X,Y,Z)\mapsto (x,y)\,,\quad x=f\frac{X}{Z}\,,\quad y=f\frac{Y}{Z}
    $$
\end{frame}

%------------------------------------------------------------
\begin{frame}
    \frametitle{Центральная проекция в однородных координатах}
    \begin{figure}
        \centering
        \includegraphics[height=0.3\textheight]{fig1.pdf}
    \end{figure}
    \begin{itemize}
        \item центральная проекция в однородных координатах
        $$
        \begin{bmatrix}
            fX\\ fY\\ Z
        \end{bmatrix}=
        \begin{bmatrix}
            f& 0& 0& 0\\
            0& f& 0& 0\\
            0& 0& 1& 0
        \end{bmatrix}
        \begin{bmatrix}
            X\\ Y\\ Z\\ 1
        \end{bmatrix}
    $$
    \item матрица камеры в случае центральной проекции
    $$
        P=\mathrm{diag}(f,f,1)\cdot[I\, |\, \mathbf{0}]\in\mathbf{R}^{3\times4}
    $$
    $$
        \mathbf{x}=P\cdot\mathbf{X}
    $$
    \end{itemize}
\end{frame}

%------------------------------------------------------------
\begin{frame}
    \frametitle{Сдвиг оптического центра}
    \begin{figure}
        \centering
        \includegraphics[height=0.3\textheight]{fig3.pdf}
    \end{figure}
    \begin{itemize}
        \item в реальных камерах начало координат не совпадает c оптическим центром
        $$
            (X,Y,Z)\mapsto\left(f\frac{X}{Z}+x_0,f\frac{Y}{Z}+y_0\right)
        $$
        \item матрица калибровки камеры
        $$
            K=
            \begin{bmatrix}
                f& 0& x_0\\
                0& f& y_0\\
                0& 0& 1
            \end{bmatrix}
        $$
        $$
            \mathbf{x}=K[I\,|\,\mathbf{0}]\cdot\mathbf{X}_\mathrm{cam}
        $$
    \end{itemize}
\end{frame}

%------------------------------------------------------------
\begin{frame}
    \frametitle{Кооординаты на изображении}
    \begin{itemize}
        \item в цифровых камерах расстояние измеряют не в метрах, а в пикселях
        $$
            K=
            \begin{bmatrix}
                m_xf& 0& x_0\\
                0& m_yf& y_0\\
                0& 0& 1
            \end{bmatrix}=
            \begin{bmatrix}
                f_x& 0& x_0\\
                0& f_y& y_0\\
                0& 0& 1
            \end{bmatrix}
        $$
        $m_x$, $m_y$ -- количество пикселей в одном метре
        \item система координат $(x,y)$ может быть косоугольной
        $$
            K=
            \begin{bmatrix}
                f_x& s& x_0\\
                0& f_y& y_0\\
                0& 0& 1
            \end{bmatrix}
        $$
        \item число степеней свободы $\mathrm{dof}(K)=5$
    \end{itemize}
\end{frame}

%------------------------------------------------------------
\begin{frame}
\frametitle{Поворот и сдвиг системы координат}
\begin{figure}
    \centering
    \includegraphics[height=0.35\textheight]{fig7.pdf}
\end{figure}

\begin{itemize}
    \item преобразование координат из $C$ в $W$
    $$
        X_w=RX_c+T
    $$
    \item матрица поворота
    $$
        R^{-1}=R^T\,,\qquad 
        \det(R)=1
    $$
    \item преобразование координат из $W$ в $C$
    $$
        X_c=R^T\left(X_w-T\right)
    $$
\end{itemize}
\end{frame}

%------------------------------------------------------------
\begin{frame}
\frametitle{Вычисление матрицы поворота}
\begin{figure}
    \centering
    \includegraphics[height=0.35\textheight]{fig6.pdf}
\end{figure}
$$
    \norm{\mathbf{r}_1}=\norm{\mathbf{r}_2}=\norm{\mathbf{r}_3}=1\,,\quad
    \mathbf{r}_1\perp\mathbf{r}_2\perp\mathbf{r}_3
$$
Матрица поворота из $xyz$ в $XYZ$
$$
    R=
    \begin{bmatrix}
        r_{11}& r_{12}& r_{13}\\
        r_{21}& r_{22}& r_{23}\\
        r_{31}& r_{32}& r_{33}
    \end{bmatrix}
    =
    \begin{bmatrix}
        \mathbf{r}_1&
        \mathbf{r}_2&
        \mathbf{r}_3
    \end{bmatrix}\,,\quad
$$
Число степеней свободы матрицы поворота:
$$
    \mathrm{dof}(R)=9-3-3=3
$$
\end{frame}

%------------------------------------------------------------
\begin{frame}
\frametitle{Экспоненциальная система координат для вращений}
\begin{figure}
    \centering
    \includegraphics[height=0.35\textheight]{fig6.pdf}
\end{figure}
По вектору вращения
$$
    \omega=
    \begin{bmatrix}
        \omega_1\\
        \omega_2\\
        \omega_3
    \end{bmatrix}
    \quad\leftrightarrow\quad
    \widehat{\omega}=
    \begin{bmatrix}
        0& -\omega_3& \omega_2\\
        \omega_3& 0& -\omega_1\\
        -\omega_2& \omega_1& 0
    \end{bmatrix}
$$
можно построить матрицу поворота (формула Родрига)
$$
    R=e^{\widehat{\omega}}=
    I+\frac{\widehat{\omega}}{\norm{\widehat{\omega}}}\sin\left(\norm{\widehat{\omega}}\right)+
    \frac{\widehat{\omega}^2}{\norm{\widehat{\omega}}^2}\left[1-\cos\left(\norm{\widehat{\omega}}\right)\right]
$$
\end{frame}

%------------------------------------------------------------
\begin{frame}
\frametitle{Экспоненциальная система координат для вращений}
\begin{figure}
    \centering
    \includegraphics[height=0.35\textheight]{fig6.pdf}
\end{figure}
По матрице поворота
$$
    R=
    \begin{bmatrix}
        r_{11}& r_{12}& r_{13}\\
        r_{21}& r_{22}& r_{23}\\
        r_{31}& r_{32}& r_{33}
    \end{bmatrix}
$$

можно вычислить вектор вращения $\widehat\omega=\log(R)$
$$
    \norm{\omega}=\cos^{-1}\left(\frac{\mathrm{trace}(R)-1}{2}\right)\,,\quad
    \omega=\frac{\norm{\omega}}{2\sin\left(\norm{\omega}\right)}
    \begin{bmatrix}
        r_{32} - r_{23} \\
        r_{13} - r_{31} \\
        r_{21} - r_{12} \\
    \end{bmatrix}
$$
\end{frame}

%------------------------------------------------------------
\begin{frame}
    \frametitle{Вращение и перемещение камеры}
    \begin{figure}
        \centering
        \includegraphics[height=0.35\textheight]{fig4.pdf}
    \end{figure}
    \begin{itemize}
        \item $\widetilde{\mathbf{C}}$, $\widetilde{\mathbf{X}}$ -- евклидовы координаты точки и центра камеры 
        в мировой системе координат
        $$
            \widetilde{\mathbf{X}}_\mathrm{cam}=
            R\cdot(\widetilde{\mathbf{X}}-\widetilde{\mathbf{C}})
        $$
        \item в однородных координатах
        $$
            \mathbf{X}_\mathrm{cam}=
            \begin{bmatrix}
                R& -R\widetilde{\mathbf{C}}\\
                0& 1
            \end{bmatrix}
            \cdot\mathbf{X}
            \quad\Rightarrow\quad
            \mathbf{x}=
            K[I\,|\,\mathbf{0}]\cdot\mathbf{X}_\mathrm{cam}=
            K\,[R\,|\,-R\widetilde{\mathbf{C}}]\cdot\mathbf{X}
        $$
        \item матрица камеры в этом случае
        $$
            P=K\,[R\,|\,\mathbf{t}]\,,\quad
            \mathbf{t}=-R\widetilde{\mathbf{C}}
        $$
    \end{itemize}
\end{frame}

%------------------------------------------------------------
\begin{frame}
    \frametitle{Внешние и внутренние параметры камеры}
    $$
            P=K\,[R\,|\,\mathbf{t}]\,,\quad
            \mathbf{t}=-R\widetilde{\mathbf{C}}
    $$
    \begin{itemize}
        \item внешние параметры камеры (extrinsics)
        $$
            R\,,\qquad \widetilde{\mathbf{C}}\,,\qquad \mathrm{dof}(R, \widetilde{\mathbf{C}}) = 3 + 3 = 6
        $$
        \item внутренние параметры камеры (intrinsics)
        $$
            K=
            \begin{bmatrix}
                f_x& s& x_0\\
                0& f_y& y_0\\
                0& 0& 1
            \end{bmatrix}
            \,,\qquad  \mathrm{dof}(K) = 5
        $$
        \item $\mathrm{dof}(P)=5+6=11$
    \end{itemize}
\end{frame}

%------------------------------------------------------------
\begin{frame}
    \frametitle{Проективная камера}
    \begin{itemize}
        \item \textbf{проективная камера}: $M\in\mathbf{R}^{3\times3}$, $\mathbf{p}_4\in\mathbf{R}^{3}$
        \[
            P=[M\,|\,\mathbf{p}_4]\in\mathbf{R}^{3\times4}\,,\qquad
            \det(M)\ne0
        \]
        \item в однородных координатах 
        $$
            \mathbf{x}=P\mathbf{x}
        $$
        $P$ определяется с точностью до множителя $\mathrm{dof}(P)=11$
        \item зная $K$, $R$ и $\mathbf{t}$ можно найти  $M$ и $\mathbf{p}_4$
        $$
            P=K\,[R\,|\,\mathbf{t}]=[KR\,|\,K\mathbf{t}]
            \quad\Rightarrow\quad
            M=KR\,,\quad
            \mathbf{p}_4=K\mathbf{t}
        $$
        \item зная $M$ и $\mathbf{p}_4$ можно вычислить $K$, $R$ и $\mathbf{t}$ c помощью матричного $\mathbf{RQ}$ разложения
    \end{itemize}
\end{frame}

%------------------------------------------------------------
\begin{frame}
    \frametitle{Анатомия проективной камеры}
    \begin{figure}
        \centering
        \includegraphics[height=0.25\textheight]{fig5.pdf}
    \end{figure}
    $$
        \mathbf{x}=P\mathbf{X}\,,\qquad
        P=[M\,|\,\mathbf{p}_4]=
        \begin{bmatrix}
            \mathbf{P}_1\\
            \mathbf{P}_2\\
            \mathbf{P}_3
        \end{bmatrix}
    $$
    \begin{itemize}
        \item $\mathbf{P}^T_2\mathbf{X}=0$ -- плоскость $y=0$, $\mathbf{P}_1$ -- плоскость $y=0$, $\mathbf{P}_3$ -- principal plane
        \item центр камеры
        $$
            P\mathbf{C}=\mathbf{0}\qquad\Rightarrow\qquad
            \mathbf{C}=
            \begin{bmatrix}
                -M^{-1}\mathbf{p}_4\\
                1
            \end{bmatrix}
        $$
        \item уравнение оптического луча, точки которого проектируются в $\mathbf{x}$
        $$
            \mathbf{X}=
            \mathbf{C}
            +\lambda
            \begin{bmatrix}
                M^{-1}\mathbf{x} \\ 
                0
            \end{bmatrix}\,,\quad
            \lambda\in\mathbf{R}
        $$
    \end{itemize}
\end{frame}

%------------------------------------------------------------
\begin{frame}
\frametitle{Поворот камеры вокруг ее центра}
\begin{figure}
    \centering
    \includegraphics[height=0.35\textheight]{fig8.pdf}
\end{figure}
Проекции для двух камер
$$
    \mathbf{x}=K[I\,|\,\mathbf{0}]\cdot\mathbf{X}
$$
$$
    \mathbf{x}'=KR[I\,|\,\mathbf{0}]\cdot\mathbf{X}=
    KRK^{-1}K[I\,|\,\mathbf{0}]\cdot\mathbf{X}=KRK^{-1}\cdot\mathbf{x}
$$
значит координаты на изображениях связаны гомографией
$$
    \mathbf{x}'=H\cdot\mathbf{x}\,,\qquad
    H=KRK^{-1}
$$
\end{frame}

%------------------------------------------------------------
\begin{frame}
\frametitle{Применение: построение панорамного изображения}
\begin{figure}
    \centering
    \includegraphics[height=0.8\textheight]{fig9.pdf}
\end{figure}
\end{frame}

%------------------------------------------------------------
\begin{frame}
\frametitle{Явление паралакса при смещении центра камеры}
\begin{figure}
    \centering
    \includegraphics[height=0.35\textheight]{fig10.pdf}
    \includegraphics[height=0.4\textheight]{fig11.pdf}
\end{figure}
\end{frame}

\end{document}
